\documentclass[11pt]{article}
\usepackage[margin=1in]{geometry}
\usepackage[T1]{fontenc}
\usepackage{lmodern}
\usepackage{amsmath,amssymb}
\usepackage{booktabs}
\usepackage{tabularx}
\usepackage{enumitem}
\usepackage[hidelinks]{hyperref}
\usepackage{xcolor}
\usepackage{microtype}
\emergencystretch=2em
\usepackage{listings}
\lstset{basicstyle=\ttfamily\small,columns=fullflexible,breaklines=true,frame=single,framerule=0.3pt}
\setlist{itemsep=2pt,topsep=3pt}

\title{CPU-LLM: A CPU-First Evidence Engine for Cited Answers,\\Deterministic Checks and Labelled Hypotheses}
\author{Amin Parva\\\small Independent researcher\\\small\url{https://github.com/insightitsGit/cpullm}}
\date{November 2026 \quad (software version 1.16.0)}

\begin{document}
\maketitle

\begin{abstract}
We describe CPU-LLM, an open-source Python library that answers questions from a user's own documents on
commodity CPUs without a generative language model. Documents are split into passages, encoded once with a frozen
sentence encoder, and stored with explicit directed facts in PostgreSQL/pgvector. A request passes through a
27-stage graph that authenticates the caller, retrieves passages through six representation channels, and
delivers only sentences or licensed constructions whose relation to the question has been checked. Every reply
carries an \texttt{answer\_status}: answered, refused for insufficient evidence, conflicting, or
unverified hypothesis. An optional strict mode delivers only literal substrings of authorized stored text.
Deterministic nodes handle rule-based deduction, exact arithmetic and interval constraints. A separate family of
nodes proposes hypotheses, always labelled \texttt{UNVERIFIED\_HYPOTHESIS}. Frontier enumerates relation chains
across a knowledge base in the manner of Swanson's literature-based discovery. Deepen then examines one hypothesis
link by link and returns supporting and contrary quotes and a test plan. We report the recorded evaluations:
a synthetic request benchmark (460/460 rows, 86/86 expected abstentions), an answerability contract (6/6
answerable, 0 unsupported answers on 9 unanswerable items), ingestion cost, and Frontier runs over three corpora
of openly licensed papers. On a holdout in which nine published findings were removed, Frontier recovered none
fully and two partially. We state what the system does not do: it does not establish that a source is true, and
its hypotheses are candidates to screen, not findings.
\end{abstract}

\section{Introduction}

Large language models produce fluent answers, but their wording is generated, and a generated claim can be
unsupported by any source \cite{ji2023hallucination}. Retrieval-augmented generation \cite{lewis2020rag} narrows
this gap by conditioning generation on retrieved passages; the final wording is still generated. Some
deployments need a different contract: an answer is a quotation or a checked construction from an authorized
source, a missing fact is reported as missing, and disagreement between sources is shown rather than resolved.
Such deployments also often lack GPUs or cannot send documents to an external service.

CPU-LLM targets that contract. Despite its name it is not a language model. It is a neuro-symbolic
\cite{garcez2020neurosymbolic} evidence engine: neural components (a frozen sentence encoder and small trained
projection heads) retrieve and route, and symbolic components (relation binders, a fact graph, deduction rules,
exact arithmetic) decide what may be said. This paper describes the whole system as of version 1.16.0, including
two recent hypothesis nodes, and collects the evaluations recorded in the repository. Our contributions are:

\begin{enumerate}
  \item an architecture in which every reply is typed by an \texttt{answer\_status} that separates source-backed
        answers, justified refusals, conflicts and unverified hypotheses (\S\ref{sec:arch}--\ref{sec:status});
  \item a strict evidence mode with an explicit, narrow meaning of ``no generated factual wording''
        (\S\ref{sec:strict});
  \item deterministic inference, arithmetic and constraint nodes with replayable proofs (\S\ref{sec:inference});
  \item Frontier and Deepen, a scored, budgeted form of literature-based discovery whose output is fenced off from
        answers (\S\ref{sec:hyp});
  \item an honest account of the evaluations and limitations (\S\ref{sec:eval}--\ref{sec:limits}).
\end{enumerate}

\section{System overview}
\label{sec:arch}

\paragraph{Storage.} Documents (text, Markdown, PDF) are split into passages with stable identifiers of the form
\texttt{prefix-p0001}. Each passage is encoded once with \texttt{all-mpnet-base-v2} \cite{reimers2019sbert} and
projected into a 1{,}024-value tensor called VectorPrism: a 16-value header, a 368-value dense core, and five
128-value slices for relational, disentangled, hyperbolic/taxonomy, identity and causal/time signals. Tensors and
text are stored in PostgreSQL with pgvector \cite{pgvector}. Only the dense and identity heads are trained in the
shipped checkpoint; the other slices are computed but not treated as enabled retrieval heads. A per-namespace fact
graph stores explicit directed facts (alias, type, location, colour, possession, lifespan, cause, contrast) with
their source sentences. Ingestion is the only encoder-heavy step; requests reuse stored vectors.

\paragraph{Request lifecycle.} A request (Python \texttt{Appliance.ask} or HTTP \texttt{POST /v1/ask}) passes a
27-stage graph. In order, the main steps are:
\begin{enumerate}
  \item authenticate the caller and resolve its source scope; an empty scope is refused before any state is used;
  \item guard the raw prompt, and guard it again after reference resolution;
  \item ask for missing criteria in subjective requests (``Which lobster is best?'') before any search;
  \item resolve pronouns against session context, or ask for an explicit name;
  \item apply conservative, corpus-backed spelling repair;
  \item encode the query and score six intent channels;
  \item retrieve dense candidates, named-entity matches and a bounded neighbourhood of the fact graph;
  \item run specialist producers (definitions, attributes, events, rankings, rules, arithmetic);
  \item bind evidence to the question (\S\ref{sec:status});
  \item audit, check grammar and render.
\end{enumerate}
Compound prompts are first split into independent cited lookups that share the caller's scope. Channel prototype
scores route a request; they never count as proof of a relation.

\paragraph{Specialist routes.} Explicit prompt prefixes or API flags send a request to a specialist node before
the ordinary graph. The routes are inference, arithmetic, CrossDomain and Constraint Check, relationship-loop
search, research, Invention, Nammade, Frontier and Deepen. A request may name only one specialist; combining two
raises an error. A classifier maps natural phrasings (``What don't we know about X?'') to a specialist only when no
existing trigger claims the prompt, so existing routes keep their behaviour.

\section{Grounded answering and the answer status}
\label{sec:status}

The binder decides whether a retrieved sentence \emph{answers} the question, not merely whether it is similar to
it. For an attribute question, a sentence qualifies only if it states that attribute of that entity, in the
question's tense, as an affirmative claim, and not under a conditional. Sharing the words ``France'' and
``capital'' is not enough, and ``France Area'' does not count as ``France''. A definition question requires a
definition, type or role construction: \emph{Mary baked a cake} cannot answer \emph{Who is Mary?}

Each reply carries an \texttt{answer\_status} and, when applicable, \texttt{answer\_evidence} (the cited passage
identifier and sentence). The statuses form four groups:
\begin{itemize}
  \item \textbf{source-backed or derived:} \texttt{answered}, \texttt{evidence\_quote}, \texttt{derived\_inference};
  \item \textbf{justified refusal:} \texttt{insufficient\_evidence}, \texttt{unsupported};
  \item \textbf{disagreement:} \texttt{conflicting}, \texttt{contradicted}, where the reply quotes each value
        and does not choose one;
  \item \textbf{unverified:} \texttt{hypotheses}, \texttt{shortage}, \texttt{research\_question},
        \texttt{deepened}.
\end{itemize}
Further statuses cover clarification, budget stops (\texttt{partial}) and guard refusals (\texttt{blocked}). An
agent or application branches on this field instead of interpreting prose.

\paragraph{Documents.} Requests such as ``Write a detailed report about X'' produce reports, articles or emails.
The text is assembled from source sentences and a small set of curated constructions. Each generated sentence is
checked by an inverse grammar: it is parsed back into the relation it claims, and that relation must match its
source proof. An optional tone (formal, friendly, and others) selects among verified constructions. General
free-form word generation is not supported. An optional external LLM fallback exists for deployments that want
it; it is off by default, and its output is labelled \texttt{external\_unverified}.

\section{Strict evidence mode}
\label{sec:strict}

Strict mode (\texttt{strict\_evidence=True}, a header, or an operator floor) intercepts a request before the
ordinary graph. It runs literal, scoped SQL evidence lookups, and no learned model writes answer wording. Delivered
evidence must be a literal substring of an authorized stored chunk. Version~2 adds bounded batches, explicit
one-hop joins, recognised paraphrase conflicts, exact scalar--unit equivalence, direct-attribute overviews and
exact arithmetic. When the evidence cannot satisfy the lookup contract, the mode returns a non-answer. It never
retries through ordinary generation. Strict mode disables every hypothesis route.

We call this ``zero hallucination'' only in a narrow sense: an execution restriction on where answer wording
comes from. It does not show that a quotation is true, current, complete or suitable for a decision. It is not a
claim of factual accuracy, universal contradiction detection or clinical validation.

\section{Deterministic inference, arithmetic and constraints}
\label{sec:inference}

\paragraph{Deduction.} ``Infer whether A is bigger than C'' or an explicit
\texttt{inference=\{subject, relation, object\}} triggers grounded deduction. Premises are source facts; the node
applies only a fixed set of application-owned rules: transitivity of size, taxonomy and part--whole relations,
membership inheritance and part--location, for example
\[
\textit{member\_of}(x,c)\wedge\textit{every\_member\_has}(c,y)\;\to\;\textit{has}(x,y).
\] It returns a replayable proof in which every premise cites a passage. The status is \texttt{derived\_inference}, or a
refusal when no proof exists. Temporal, inductive and abductive analyses run over explicitly selected JSON
evidence records, and their outputs are labelled by kind. One prepared deduction is shared by the research,
Nammade and CrossDomain routes in a request.

\paragraph{Arithmetic and constraints.} ``Calculate 17 times 24'' is evaluated exactly (decimal and rational
arithmetic) rather than generated. Constraint Check intersects numeric intervals stated in sources. CrossDomain joins facts
across knowledge-base prefixes. A finance module checks ledger covenants. Relationship-loop search
(\texttt{cycle\_search}) enumerates bounded cycles in the fact graph under explicit round, candidate and time
budgets.

\section{Hypothesis nodes}
\label{sec:hyp}

The remaining nodes produce candidates rather than answers. Every result is labelled
\texttt{UNVERIFIED\_HYPOTHESIS}, comes with \texttt{promotion: blocked}, is never written back to the index, and carries a standard disclosure.

\paragraph{Research, Invention, Nammade and the abductive bridge.} ``Research X'' runs Innovation (taxonomy-driven
cross-passage invention), source hypothesis generation and a Nammade scan in parallel on one snapshot, and lists
each reading as a separate option with the reason it abstains or proposes. A Nammade scan quotes authorized
passages admitted by the dense core (cosine~$\ge 0.62$) and reports volumetric labels and polarity bridges between
words that share a span; those labels are explicitly not evidence. The abductive bridge emits one schema-bound
hypothesis when two passages give interval-disjoint bounds on one observable and differ on a regime variable.

\paragraph{Frontier.} Frontier implements a scored, budgeted variant of Swanson's ABC model of literature-based
discovery \cite{swanson1986fish,swanson1997arrowsmith,henry2017lbd}: if one document links A to B and another links
B to C, the A--C connection is a candidate no single document states. Frontier extracts relation edges
(cause, increase, decrease, constrain, and others, with hedging and negation) from every authorized passage.
It uses a spaCy dependency parse when a model is installed \cite{honnibal2020spacy} and a regex fallback otherwise.
Around the topic it streams three candidate shapes: chains $A\to B\to C$, convergent pairs (two causes of one
outcome) and divergent pairs (one cause, two outcomes). The search is bounded by rounds, candidates, edges, memory
and time (default 60\,s); a stopped run is flagged \texttt{incomplete}. Candidates pass three filters in order:
\begin{itemize}
  \item \emph{abductive:} the edges compose into a coherent explanation;
  \item \emph{deductive:} each edge has an exact source span, and no source states its negation or its opposite;
  \item \emph{pre-thinker:} rejects malformed, circular, single-source and restated candidates.
\end{itemize}
A candidate stated as a whole by any one passage is not new and is dropped.

Each survivor receives three scores in $[0,1]$. An edge's support is $0.80$ for an assertive source and
$0.55$ for a hedged one. It gains $0.10$ per additional supporting document (at most $0.20$) and $0.05$ for a
quantitative statement, capped at 1. A path's support is the minimum over its edges, and
\begin{equation}
  \textit{source\_support} = 1 - \prod_{i=1}^{\min(k,5)} (1 - s_i),
\end{equation}
a noisy-or over up to five supporting paths. Relevance combines composition certainty $c$, quote fit $q$ and
topic-term fit $t$, scaled by term specificity:
\begin{equation}
  \textit{relevance} = (0.5c + 0.3q + 0.2t)\,(0.6 + 0.4\sigma), \qquad
  \sigma = \frac{1}{|T|}\sum_{w\in T}\Bigl(1 - \frac{\log(1+\mathrm{df}_w)}{\log(1+N)}\Bigr),
\end{equation}
where $\mathrm{df}_w$ is the number of passages containing term $w$ and $N$ the number of passages. Without the
specificity factor, generic hubs such as ``performance'' dominated the ranking. Novelty is
\begin{equation}
  \textit{novelty} = 0.5(1 - o) + 0.3\min\!\Bigl(1, \frac{D-1}{2}\Bigr) + 0.2\,\textit{sep},
\end{equation}
with $o$ the largest fraction of the hypothesis's content words found in any one premise quote, $D$ the number
of distinct source documents, and $\textit{sep}=0.5$ when some sentence names both endpoints (without a relation cue),
$0.75$ when they share only a passage, and $1$ when they never co-occur. Entries are ranked by support, then relevance, then novelty.
No term may appear in more than three of the leading entries. Frontier returns 10 to 50 entries; when fewer survive,
it reports a \texttt{shortage} rather than padding the list. Each entry carries its citations and a falsification
test. When an encoder is available, a link between two terms that share only a head noun must also pass a
cosine check ($\ge 0.6$).

\paragraph{Deepen.} Deepen takes one hypothesis (a Frontier entry, free text, or ``Deepen idea 3'' in a session)
and examines it link by link. For each link it searches the whole authorized graph, not only the passages that
produced the hypothesis, for supporting and contrary statements, and assigns a status: \texttt{supported},
\texttt{contested}, \texttt{counter\_only}, \texttt{weak} or \texttt{missing}. The overall verdict is one of
\texttt{all\_links\_supported}, \texttt{partly\_supported}, \texttt{contested}, \texttt{unsupported} or
\texttt{incomplete} (a budget stop). It also reports:
\begin{itemize}
  \item source predictions that name at least two of the hypothesis's terms;
  \item existing experimental setups that already measure one of its terms;
  \item confounders, i.e.\ other stated causes of the outcome, with their causal role (outcome predictor, shared
        cause or downstream) but never adjusted for automatically \cite{pearl2009causality};
  \item a test plan (mediation, $2\times2$ factorial or common cause) stating what to vary, measure, block and
        control, and when the hypothesis fails.
\end{itemize}
Version 1.16.0 adds a composition assessment separate from per-link support. Composition certainty is 1.0 for an
exact shared term and lower for weaker term matches (0.95, 0.85, 0.7, 0.6), multiplied by $0.8$ per head-noun
bridge (or by the cosine when an encoder is present) and by $0.9$ per edge beyond two. The release also adds
warnings when a bridged term changes modifier or meaning (``low weight decay'' vs.\ ``increasing weight decay''),
user-supplied definitions and protocol, a readiness flag and an optional novelty review against the knowledge
base. Frontier runs Deepen on its first five entries by default and flags contested ones with the contrary quote.

Because every Frontier link has at least its own premise, \texttt{all\_links\_supported} is nearly guaranteed for
Frontier output and is weak evidence. The informative signals are \texttt{contested} and \texttt{missing}.

\section{Evaluation}
\label{sec:eval}

All numbers below are recorded in the repository with their scripts and revisions. They measure stated contracts
on stated fixtures; none is a general accuracy figure.

\subsection{Grounded answering}

\paragraph{Answerability contract.} Seventeen frozen fixtures were scored three ways: six answerable, nine
unanswerable, and two conflicts. Table~\ref{tab:answer} gives the results. Baseline~A labels the first retrieved
passage as a passage. Baseline~B speaks the first passage as the answer. The fixtures were written with the binder,
so this is a contract test rather than a held-out corpus.

\begin{table}[h]
\centering\footnotesize
\begin{tabular}{lcccc}
\toprule
 & Answerable & Unsupported & Conflicts & Exact-span\\
 & covered (of 6) & answers (of 9) & shown (of 2) & citations\\
\midrule
Baseline A & 0 & 0 & 0 & 0\\
Baseline B & 5 & 8 & 0 & 5\\
CPU-LLM binder & 6 & 0 & 2 & 6\\
\bottomrule
\end{tabular}
\caption{Answerability contract (fixtures in \texttt{tests/test\_answerability.py}).}
\label{tab:answer}
\end{table}

\paragraph{Request benchmark.} A synthetic request suite exercises citation, access control, clarification,
abstention and budgets. At revision \texttt{7adf319}, 460/460 rows passed, including 278/278 expected answers and
86/86 expected abstentions, and the scorer flagged no false assertions. An independent harbour fixture first
scored 20/24. Both misses involved the colour word \emph{ochre}. After a vocabulary fix the fixture replayed at
24/24, and a typo replay passed 32/32. The replay figures measure regression coverage after a fix, not first-run
accuracy. Our own validity review notes that the aggregate mixes response depth and adversarial input handling
with evidence lookup, and should be read as a regression metric.

\paragraph{Other receipts.} The stylistic document node passed 49 contract scenarios per arm (rules and spaCy). A
healthcare demonstration on fictional data produced 12/12 authored outcomes; this is not clinical validation. The source test suite at 1.16.0 gives 1{,}343 passed, 56 skipped (mostly opt-in database
integration tests) and 3 expected failures, with 3 failures that predate this release.

\subsection{Cost on a CPU}

A 162-page PDF (MIT 18.04 lecture notes) produced 265 chunks. With one encoder worker and four Torch threads,
ingestion took 85.7\,s and peaked at 3.00\,GiB. With four workers, it took 113.7\,s and 6.53\,GiB. Both runs
stored identical chunks and vectors. On a seven-question latency check after start-up, the median request took
167\,ms (range 112--668\,ms), with 11.5\,s start-up. That check also recorded failures: ``What is the capital of
France?'' returned unrelated stored sentences at that revision. The answerability binder above was the
response. Frontier on a 1{,}054-chunk knowledge base took about 3.5\,s per prompt at 2.98\,GB resident memory.

\subsection{Hypothesis nodes}

We built three knowledge bases from openly licensed arXiv papers (fixtures in the repository):
\begin{itemize}
  \item AI bottlenecks: 64 papers, 32 in full text (CC BY/BY-SA), 1{,}054 chunks;
  \item deep-learning generalization: 106 papers, 60 in full text, 1{,}652 chunks;
  \item consciousness and philosophy of mind: 153 papers, 48 in full text, 1{,}241 chunks.
\end{itemize}

\paragraph{Holdout rediscovery.} In the AI-bottlenecks base we selected nine published findings and removed every
sentence that contained the terms of any finding: 767 sentences. We then asked whether Frontier would propose them.
An automatic substring matcher reported 3 full and 7 partial matches. Hand adjudication found that most of these
were spurious, and the corrected result is \textbf{0 of 9 full and 2 of 9 partial}. One partial is close to text
that remained in another paper and is weak evidence of discovery. Three repeated runs produced identical output.

\paragraph{Deepen checks.} On the deep-learning base, Deepen checked 70 leading Frontier entries: 65 had all links
supported, 4 were partly supported and 1 was contested. The contested entry was ``low weight decay may increase
generalization through the grokking phase''; a corpus paper \cite{clauw2024grokking} reports the opposite effect.
An early version found 69/70 supported because it searched only the passages that produced the hypothesis, which
hides contrary evidence; it now searches the whole graph, and a regression test covers this. Separately, on seven
hypotheses Deepen found two contested. The second was a claim about implicit bias, negated by a ``neither \ldots{}
alone'' sentence in \cite{matveev2025biases}. With the 1.16.0 composition review, four of five fully supported
hypotheses were marked \texttt{requires\_review} because a bridge changed a term's meaning, and the novelty review
found prior overlap for one of seven (\cite{mayama2025iit}). Frontier's output on 15 prompts was unchanged by the
review layer.

\paragraph{Qualitative observations.} Causal empirical prose (systems and ML papers) gave the most usable chains.
Philosophy gave the fewest: its arguments are rarely stated as directed causal relations. Several entries read as
plausible research leads to us, but we did not run a blind expert rating like that of Si et al.\ \cite{si2024ideas}.
We therefore claim only that a minority of entries are worth a researcher's attention, not that they are correct
or novel to the field.

\section{Related work}

Retrieval-augmented generation \cite{lewis2020rag} and the hallucination literature \cite{ji2023hallucination}
motivate grounding; CPU-LLM removes generation from the answer path rather than conditioning it. Neuro-symbolic
AI \cite{garcez2020neurosymbolic} combines learned perception with symbolic reasoning; here the learned parts
retrieve and route while symbolic parts decide what may be said. Literature-based discovery
\cite{swanson1986fish,swanson1997arrowsmith,henry2017lbd} is the closest precedent for Frontier. Our additions
are typed and hedged relation edges, explicit scoring, budgets, deductive and contrary-evidence filtering, and a
link-level audit (Deepen). LLM-driven discovery systems such as FunSearch \cite{romeraparedes2024funsearch} pair
generation with an automatic evaluator; Frontier has no generator, and its evaluator is source evidence rather
than an executable objective. Human studies of LLM research ideation \cite{si2024ideas} indicate the kind of
evaluation our hypothesis nodes still need.

\section{Limitations}
\label{sec:limits}

\begin{itemize}
  \item \textbf{Source truth.} A cited quotation can be false, outdated or out of context. The system checks that
        a sentence states the asked relation, not that the world agrees.
  \item \textbf{Coverage.} The binder handles recognised question shapes; other forms abstain. Word overlap and
        parse patterns are not semantic entailment. Free-form generation is not supported.
  \item \textbf{Trained components.} Only the dense and identity heads are trained; four of six channels are
        symbolic or uncalibrated.
  \item \textbf{Evaluation.} Most benchmarks are synthetic or were written alongside the code. The holdout result
        is small (nine findings) and was hand adjudicated by the developer. There is no blind expert rating of
        hypotheses and no comparison against an LLM baseline on the same corpora.
  \item \textbf{Hypotheses.} Frontier depends on how explicitly a field states directed relations. Its scores
        are heuristics with hand-chosen constants, not calibrated probabilities.
  \item \textbf{Deployment.} Ingestion requires PostgreSQL with pgvector 0.7 or later and about 3\,GB of memory.
        The fact graph is scanned rather than indexed.
\end{itemize}

\section{Availability and reproducibility}

CPU-LLM is available on PyPI as \texttt{cpu-llm} and at \url{https://github.com/insightitsGit/cpullm}
\cite{cpullm2026}. This paper describes version 1.16.0. The repository contains:
\begin{itemize}
  \item the three corpora with their source lists, licences and holdout definitions (\texttt{tests/fixtures});
  \item the fetch, build, run and evaluation scripts and raw outputs (\texttt{docs/reviews});
  \item benchmark and answerability records under \texttt{docs/};
  \item an OpenAPI contract and function-calling tool definitions for AI agents (\texttt{docs/}).
\end{itemize}
Frontier and Deepen are deterministic for a fixed corpus, parser model and configuration.

\section*{Acknowledgements and AI assistance}
The system design, including the Frontier and Deepen nodes, is the author's. Parts of the implementation, tests,
evaluation scripts and the drafting of this manuscript were done with an AI coding assistant (Claude, Anthropic)
under the author's direction. The author reviewed the content and takes full responsibility for it. The corpus
papers are by their respective authors and are used under their open licences.

\bibliographystyle{plain}
\bibliography{references}

\end{document}
